\subsection{Explainable Fault Localization}
\label{sec:explainable-fl}

Given a bug report, our fault localization procedure follows a \emph{progressive refinement} strategy over the Concurrent Knowledge Base (CKB).
Rather than exposing large code fragments or entire program graphs to an LLM, we localize concurrency faults by performing hierarchical retrieval and reasoning over structured, DSL-encoded concurrency knowledge.

The hierarchical design is motivated by the structural properties of large software systems.
Well-engineered systems are typically organized with high cohesion and low coupling: components differ significantly in functionality, packages encapsulate closely related behaviors, and thread entry points define distinct concurrent execution contexts.
As a result, concurrency faults are usually confined to a small region of the system and rarely span arbitrary components or unrelated execution contexts.
This observation motivates a coarse-to-fine localization strategy that progressively narrows the search space along natural architectural boundaries.

\paragraph{Hierarchical Localization Workflow.}
Starting from the bug report, the LLM incrementally refines the localization scope by querying the CKB in the following order:
(i) selecting the most relevant component,
(ii) selecting a package within the chosen component,
(iii) localizing candidate thread entry points within the package, and
(iv) selecting a pair of interacting thread entries.
At each stage, only a small and relevant slice of the CKB is exposed to the LLM (e.g., a list of components, packages, or thread entries), ensuring that the reasoning context remains bounded and focused.

Once a thread-entry pair is selected, the framework retrieves all interaction-level records associated with the pair.
These records are represented using the interaction DSL and summarize shared resources, access actions, and synchronization context when available.
The interaction set provides a compact, concurrency-focused context that identifies the candidate methods involved in the interaction.

\paragraph{LLM-based Interaction Ranking and Explanation.}
Given the interaction DSL elements, the LLM performs interaction-level reasoning to assess which methods are most likely responsible for the fault.
To support this reasoning, the LLM may selectively query method-level summaries and, when necessary, corresponding code snippets through predefined retrieval functions.
Based on the combined interaction context and method-level information, the LLM produces a ranked list of suspicious methods.

Finally, the LLM generates an explanation grounded in explicit concurrency context.
The explanation links the ranked methods to their accessed shared variables and the corresponding thread-entry pair, describing how concurrent execution and shared-resource access may lead to the observed faulty behavior.

\begin{figure}[H] \scriptsize
\centering 
\vspace*{3pt}
\begin{tcolorbox}[colback=gray!5,colframe=black!30]

You are an expert in concurrent fault localization. Given a bug report and access to a Concurrent Knowledge Base (CKB), iteratively retrieve relevant concurrency information through hierarchical reasoning and identify the most suspicious methods related to the reported fault.

\vspace{0.3em}

\textbf{Bug Report:} \verb!{{BUG_REPORT}}!

\# Component-Level Retrieval

First, call: \textit{$list\_components()$}. Use function calls to retrieve concurrency-related knowledge from the CKB, then analyze the returned components and identify the most relevant component associated with the bug report.

\begin{verbatim}
{{function_call}}
\end{verbatim}
\textbf{Output Format:} component\_name: component
\vspace{0.3em}

\# Package-Level Retrieval

First, call: \textit{$list\_packages\_under\_component(component)$}. Use function calls to retrieve concurrency-related knowledge from the CKB, then analyze the returned packages and identify the most relevant package associated with the bug report. 

\begin{verbatim}
{{function_call}}
\end{verbatim}
\textbf{Output Format:} package\_name: package
\vspace{0.3em}

\# Thread Entry and Method-Level Retrieval

First, call: \textit{$list\_entries\_under\_package(package)$}. Use function calls to retrieve concurrency-related knowledge from the CKB, then analyze the returned thread entry points and identify the most suspicious entry pair related to the reported fault. 

Select the most suspicious thread-entry pair.

\begin{verbatim}
{{function_call}}
\end{verbatim}
\textbf{Output Format:} entry\_pair: (entry1, entry2)
\vspace{0.3em}

\# Interaction DSL Reasoning

First, retrieve the interaction DSL: $get\_interaction\_dsl(entry_pair)$. Use the DSL to analyze shared resources, access actions, synchronization context, and cross-thread interactions. If more evidence is needed, continue using function calls to inspect relevant method summaries or source code:
\begin{verbatim}
{{function_call}}
\end{verbatim}

After sufficient evidence has been collected, apply the reasoning rules to identify suspicious methods involved in the concurrent interaction. Then return a ranked list of suspicious methods and explain how the interaction may lead to the reported fault.

\textbf{Output Format:}

\begin{verbatim}
Rank 1: [Method]: ... [Shared Resource]: ... [Thread Entries]: ... [Explanation]: ...
\end{verbatim}
\end{tcolorbox}
\vspace{-0.4em}
\caption{Prompt Template} \label{template}
\vspace{-0.15in}
\end{figure}


\paragraph{Prompt Template.} 
ConFL follows the function-calling paradigm common to recent LLM-based FL approaches (e.g., AutoFL, FlexFL). Figure~\ref{template} illustrates the complete prompt template. At each hierarchical retrieval stage, the LLM is provided with the bug report and a typed set of retrieval functions (listed in Table~\ref{tab:function_calls}); it selects one function to invoke, observes the structured response, and progresses to the next stage. This constrains the LLM to operate within the CKB schema rather than generating free-form responses. At the interaction-ranking stage, each interaction DSL element is serialized as a labeled block listing the entry pair, the set of shared methods, and the per-method shared-variable accesses with their guard locks, so that the LLM can reason directly over concurrency-relevant context without traversing raw code. We use deterministic decoding (temperature = 0) and the complete prompt templates are released in the artifact.



 